\subsection{Related work}

Prior music systems connect listener information with song coordinates in
several ways. \emph{Feel the Moosic} lets the listener choose a position in a
simplified affect space and requests tracks from corresponding Spotify
Valence--Energy ranges~\cite{helmholz2019}. Janssen et al.\ learned
listener-specific relations between music and physiological response and used
them to steer an affective music player toward a goal~\cite{janssen2012}.
Ayata et al.\ estimated Valence and Arousal from wearable
galvanic-skin-response and photoplethysmography measurements~\cite{ayata2018},
while Tran et al.\ matched a facial estimate with nearby tracks in Spotify
feature space~\cite{tran2023}. These systems establish several routes from
explicit or inferred context to music.

Two issues remain. First, classification accuracy does not establish that a
recommended track improves the listening experience. Second, the path from
sensor observation to selection can hide consequential assumptions.
Explanations may support transparency, trust, effectiveness, or
decision-making~\cite{nunes2017}, while user experience also depends on effort,
privacy, and context~\cite{knijnenburg2012}. A reviewable policy should expose
what was observed, how it was summarized, and why one song became eligible.

\subsection{Research gap, questions, and contributions}

Among the systems reviewed here, few expose a baseline-relative path from
multimodal observations to a reviewable next-song decision. It also remains
unclear whether that policy changes listener judgments when catalog, exposure,
and rating procedure are held constant. Vibe Shuffle addresses this gap with a
personal reference, an event-driven listening trajectory, a fixed
region-and-distance rule, and a Random comparator drawn from the same frozen
catalog.

\paragraph{RQ1.}
How can within-listener changes in facial and cardiac observations be
transformed into an inspectable next-song decision without assigning absolute
emotion labels?

\paragraph{RQ2.}
How does Vibe Shuffle compare with Random song selection in participants'
ratings of mood fit and liking?

The paper makes three concrete contributions. First, it defines a
baseline-relative formulation of multimodal music adaptation. Second, it
specifies the path from personal references and listening trajectories to an
eligible-song decision. Third, it reports a matched exploratory comparison
that quantifies the direction and uncertainty of mood-fit and liking
differences. These contributions define a testable policy and its current
evidence boundary.

In the exploratory comparison, mood-fit ratings were 0.56 points higher under
Vibe Shuffle on average, but the confidence interval crossed zero. Liking
differed by 0.09 points and was similarly uncertain. These estimates motivate a
better instrumented test of the policy; they do not show that the system
improves listening.
